\documentclass{vvsu}

\vvsuyear{2026}

%%%%%%%%%%%%%%%%%%%

\usepackage{graphicx} % для изображений
\usepackage{tabularray} % для таблиц
\usepackage{siunitx} % для обозначений (процент, градус)
\usepackage{csquotes} % для корректной работы с кавычками
\usepackage{listings} % для листингов кода

% Список путей, где будут искаться изображения и файлы
\graphicspath{{images/}}

% Автор документа
\author{А. Е. Филатов}

% Настройка стилей для листингов кода
\input{listing_styles.tex}

%%%%%%%%%%%%%%%%%%%

\begin{document}

% Шапка
\vvsuhead{\linespread{1}\selectfont{}МИНОБРНАУКИ РОССИИ\\
\vspace{10pt}Федеральное государственное бюджетное образовательное учреждение\\
высшего образования\\
\fontsize{13}{13}\selectfont{}<<ВЛАДИВОСТОКСКИЙ ГОСУДАРСТВЕННЫЙ УНИВЕРСИТЕТ>>\\
(ФГБОУ ВО <<ВВГУ>>)\\
\vspace{10pt}\fontsize{12}{12}\selectfont{}ИНСТИТУТ ИНФОРМАЦИОННЫХ ТЕХНОЛОГИЙ И АНАЛИЗА ДАННЫХ\\
КАФЕДРА ИНФОРМАЦИОННЫХ ТЕХНОЛОГИЙ И СИСТЕМ}

% Название отчета
\title{Отчет\\по лабораторной работе №6}
\subtitle{по дисциплине\\<<Информатика и программирование>>}

% Участники работы
\member{Студент\\ гр. БИН-25-3}{А. Е. Филатов}
\member{Ассистент\\ преподавателя}{М.В. Водяницкий}

% Вывод титульника
\maketitle

% Задание
\begin{addition}{Задание}
  Выполнить задания и оформить отчет по стандартам ВВГУ.

  \textit{\textbf{Задание 1.}}
  Написать функцию, которая конвертирует время из одной величины в другую.
  
  На вход подается:
  \begin{vvsu_itemize}
    \item число (величина времени)
    \item исходная единица измерения
    \item единица измерения, в которую нужно перевести
  \end{vvsu_itemize}
  Функция должна вернуть конвертированное значение

  Примеры (формат ввода/вывода можно выбрать свой, если нет строгих требований):
  \begin{vvsu_table}{}
  \label{table:space_terms}
  \begin{tblr}{
    width=1\linewidth,
    colspec={lX},
    hlines, vlines
  }
  Вход	   & Выход \\
  4h m	   & 240m \\
  30m h    & 0.5h \\
  12s h	   & 0.03h \\
  \end{tblr}
  \end{vvsu_table}

  \textit{\textbf{Задание 2.}}  
  Пользователь делает вклад в банке в размере a рублей сроком на n лет

  Процент по вкладу зависит от суммы и срока

  Зависимость от суммы:
  \begin{vvsu_itemize}
    \item каждые 10 000 рублей увеличивают ставку на 0.3 процента
    \item но суммарное увеличение не может превышать 5 процента
    \item минимальный вклад - 30 000 рублей
  \end{vvsu_itemize}

  Зависимость от срока:
  \begin{vvsu_itemize}
    \item первые 3 года - 3 процента
    \item от 4 до 6 лет - 5 процента
    \item более 6 лет - 2 процента
  \end{vvsu_itemize}
  Необходимо написать функцию, которая рассчитывает прибыль пользователя без учета первоначально вложенной суммы

  Используется сложный процент: каждый год процент начисляется на текущую сумму вклада

  На вход подаются: сумма вклада и количество лет. Результат: сумма прибыли (не весь вклад, а только заработанные проценты)

  Примеры:
    \begin{vvsu_table}{}
  \label{table:space_terms}
  \begin{tblr}{
    width=1\linewidth,
    colspec={lX},
    hlines, vlines
  }
  Вход	   & Выход \\
  30000 3	 & 3648.67 \\
  100000 5 & 38920.10 \\
  200000 8 & 183925.42 \\
  \end{tblr}
  \end{vvsu_table}

  \textit{\textbf{Задание 3.}}  
  Написать функцию для вывода всех простых чисел в заданном диапазоне. Нужно учитывать некорректные данные (например, начало больше конца или диапазон без простых чисел)

На вход подаются два числа: начало и конец диапазона (включительно). На выходе - список всех простых чисел или сообщение об ошибке

Примеры:
    \begin{vvsu_table}{}
  \label{table:space_terms}
  \begin{tblr}{
    width=1\linewidth,
    colspec={lX},
    hlines, vlines
  }
  Вход	   & Выход \\
  1 10	 & 2 3 5 7 \\
  15 120 & 17 19 23 29 31 37 41 43 47 53 59 61 67 71 73 79 83 89 97 101 103 107 109 113 \\
  0 1    & Error! \\
  \end{tblr}
  \end{vvsu_table}
  (Формат вывода списка простых чисел может быть любым удобным: в строку через пробел, в несколько строк и т.п.)

  \textit{\textbf{Задание 4.}}  
  Реализовать функцию сложения двух матриц

  При сложении двух матриц получается новая матрица того же размера, где каждый элемент - это сумма элементов с тем же индексом из двух исходных матриц

  Ограничения:
  \begin{vvsu_itemize}
    \item складывать можно только матрицы одинакового размера
    \item размер матрицы должен быть строго больше 2 (например, 3×3, 4×4 и т.д.)
    \item при нарушении условий нужно вывести сообщение об ошибке
  \end{vvsu_itemize}

  На вход подаются:
  \begin{vvsu_list}
    \item размер матрицы n (для квадратной матрицы n × n);
    \item элементы первой матрицы (по строкам, через пробел);
    \item элементы второй матрицы в таком же формате.
  \end{vvsu_list}

  Результат - новая матрица (в том же формате), либо сообщение об ошибке

  \textit{\textbf{Задание 5.}}  
  Написать функцию, которая определяет, является ли строка палиндромом

  Палиндром - это строка, которая читается одинаково слева направо и справа налево (обычно без учета пробелов, регистра и знаков препинания - эти правила нужно явно задать в своей реализации)

На вход подается строка. На выходе:
   \begin{vvsu_itemize}
    \item Да, если это палиндром
    \item Нет, если это не палиндром
    \end{vvsu_itemize}

Примеры:
    \begin{vvsu_table}{}
  \label{table:space_terms}
  \begin{tblr}{
    width=1\linewidth,
    colspec={lX},
    hlines, vlines
  }
  Вход	   & Выход \\
  А роза упала на лапу Азора	 & Да \\
  Borrow or rob & Да \\ 
  Алфавитный порядок   & Нет \\
  \end{tblr}
  \end{vvsu_table}

\textit{\textbf{Оформление отчета}}  

Отчет оформляется строго по СТО.

Не забудьте добавить страницу "Задание" с копией содержимого этого файла (с правильным оформлением списков и т.д.)

В отчете должно быть объяснено как работает ваша программа (каждое отдельное задание)

\end{addition}

% Содержание
\toc

% Глава - Выполнение работы
\section{Выполнение работы}
\label{sec:execution}

% Подглава - Задание 1
\subsection{Задание 1}
  Сначала создадим функцию convert, конвентирум с одной величины в другую, а после выводим результат. На рисунке \ref{fig:../code/6.1.py} представлен код программы.

\begin{vvsu_figure}{Листинг программы для задания 1}{fig:../code/6.1.py}
  \begin{minipage}{.75\textwidth}
    \lstinputlisting[language=Python,basicstyle=\fontsize{10}{10}\linespread{1}\selectfont\ttfamily]{../code/6.1.py}
  \end{minipage}
\end{vvsu_figure}

\begin{vvsu_list}
  \item создаем функцию convert;
  \item создаем переменнут time для ввода через пробел данные;
  \item переводим все 3 функции в нужные нам для роботы типы данных
  \item проверяем по условиям и возвращаем подсчитанное значение;
  \item если чтото иное введено, то выводим "некорректные данные";
  \item выводим функцию с результатом.
\end{vvsu_list}

% Подглава - Задание 2
\subsection{Задание 2}
 Создаем функцию deposit, вводи значения, а дальше считаем процент по вкладу по условию здания, выводим рузультат с вызовом функции. На рисунке  \ref{fig:../code/6.2.py} представлен код программы.

\begin{vvsu_figure}{Листинг программы для задания 2}{fig:../code/6.2.py}
  \begin{minipage}{.75\textwidth}
    \lstinputlisting[language=Python,basicstyle=\fontsize{10}{10}\linespread{1}\selectfont\ttfamily]{../code/6.2.py}
  \end{minipage}
\end{vvsu_figure}

\begin{vvsu_list}
  \item создаем функцию deposit;
  \item просим на ввод от пользователя сумму и срок вклада через пробел;
  \item разделяем каждое значеение и преобразуем их в нужный тип данных;
  \item проверяем условие и от условия задания считаем процент по вкладу;
  \item возвращаем прибыль;
  \item выводим функию с подсчитанным результатом.
\end{vvsu_list}

% Подглава - Задание 3
\subsection{Задание 3}
  Создаем функцию prime num, вводи промежуток чисел через пробле, ищем простые числа и добавляем их в массив, выводим рузультат с вызовом функции. На рисунке \ref{fig:../code/6.3.py} представлен код программы.

\begin{vvsu_figure}{Листинг программы для задания 3}{fig:../code/6.3.py}
  \begin{minipage}{.75\textwidth}
    \lstinputlisting[language=Python,basicstyle=\fontsize{10}{10}\linespread{1}\selectfont\ttfamily]{../code/6.3.py}
  \end{minipage}
\end{vvsu_figure}

\begin{vvsu_list}
  \item создаем функцию prime num;
  \item просим у пользователя на ввод промежуток чисел через пробел и разделяем их;
  \item проебазуем их их в нужный тип данных;
  \item создаем массив;
  \item проверяем по условию каждое число в промежутке заданных чисел и ищем наличие ошибок, несоответсвией;
  \item если все хорошо, добавляем число в массив
  \item возвращаем объединнные все строки в одну и по разделителю пробел;
  \item выводим функцию с подсчитанным результатом.
\end{vvsu_list}

% Подглава - Задание 4
\subsection{Задание 4}
  Создаем функцию которая проверяет состоит ли кортеж только из чисел, если да, то сортируем числа от еньшего к большему, если неь, то кортеж остается неизменным. На рисунке  \ref{fig:../code/6.4.py} представлен код решения.

\begin{vvsu_figure}{Листинг программы для задания 4}{fig:../code/6.4.py}
  \begin{minipage}{.75\textwidth}
    \lstinputlisting[language=Python,basicstyle=\fontsize{10}{10}\linespread{1}\selectfont\ttfamily]{../code/6.4.py}
  \end{minipage}
\end{vvsu_figure}

\begin{vvsu_list}
  \item объявляем функцию sort с параметром х;
  \item проверка, все ли элементы в x - числа (int или float);
  \item если возвращает значение True, то сортируем кортеж и возвращаем отсортированный кортеж;
  \item если возвращает значение False, то возвращаем кортеж без изменений;
  \item проверка работы функции значением True;
  \item проверка работы функции с значением False.
\end{vvsu_list}

% Подглава - Задание 5
\subsection{Задание 5}
  Создае функцию которая получает на ввод словарь и выведет пользователю товар с максимальной и минимальной ценой. На рисунке  \ref{fig:../code/6.5.py} представлен код программы.

\begin{vvsu_figure}{Листинг программы для задания 5}{fig:../code/6.5.py}
  \begin{minipage}{.75\textwidth}
    \lstinputlisting[language=Python,basicstyle=\fontsize{10}{10}\linespread{1}\selectfont\ttfamily]{../code/6.5.py}
  \end{minipage}
\end{vvsu_figure}

\begin{vvsu_list}
  \item импортируем встроенный модуль для работы с регулярными выражениями;
  \item создаем функицю palindrome;
  \item прееводим все символы переменной к нижнему регистру;
  \item удаляем с прееменной все символы, кроме букв и цифр;
  \item проверяем, если переменная равна своей же обратной версии, то возвращаем "Да", иначе "Нет";
  \item ипользуе функцию с разными примерами.
\end{vvsu_list}

\end{document}